\section{Как получена формула (25) Сеничева и что такое $F_v$}
\label{sec:f25}

В электростатическом дефлекторе частица на смещённой орбите проходит другой потенциал. Её задержка за оборот складывается из удлинения пути и из вклада изменённой скорости на этой траектории. В~\cite{Senichev2013} вторая часть заменена одним множителем $F_v$. Ниже формула (25) той работы выписана по странице proceedings, множитель сопоставлен с текстом IPAC'11 и с рядом раздела~\ref{sec:es}.

\subsection{Что напечатано}
\label{sec:f25-print}

Магнитный период в~\cite{Senichev2013} собирается из длины и скорости:
\begin{equation}
\label{eq:f25-5}
\frac{\Delta T_{\mathrm{rev}}}{T_{\mathrm{rev}}}
=
\frac{\Delta C}{C}
-
\frac{\Delta v}{v_s}
-
\frac{\Delta C}{C}\cdot\frac{\Delta v}{v_s}
+
\left(\frac{\Delta v}{v_s}\right)^2.
\end{equation}
При $\Delta C/C=\alpha_0\delta+\alpha_1\delta^2$ и при $\Delta v/v_s=\delta/\gamma^2$ без члена $\delta^2$ правая часть~\eqref{eq:f25-5} даёт две скобки ниже. Удлинение пути от бетатронных колебаний работа дописывает отдельно, и вместе получается формула (6) той статьи:
\begin{equation}
\label{eq:f25-6}
\frac{\Delta T_{\mathrm{rev}}}{T_{\mathrm{rev}}}
=
\left(\alpha_0-\frac{1}{\gamma^2}\right)\delta
+
\left(\alpha_1-\frac{\alpha_0}{\gamma^2}+\frac{1}{\gamma^4}\right)\delta^2
+
\left(\frac{\Delta L}{L}\right)_{\!\beta}.
\end{equation}
Скобка при $\delta^2$ --- та же, что в формуле (5) работы: слагаемого $3\beta_s^2/(2\gamma_s^2)$ в ней нет. В обозначениях~\eqref{eq:set-tau} полный магнитный коэффициент раздела~\ref{sec:mag} равен
\begin{equation}
\label{eq:f25-mag}
\eta_1
=
\alpha_1-\frac{\alpha_0}{\gamma_s^2}+\frac{1}{\gamma_s^4}+\frac{3\beta_s^2}{2\gamma_s^2}.
\end{equation}
Сверх скобки~\eqref{eq:f25-6} в~\eqref{eq:f25-mag} стоит квадратичный член свободной скорости. Вне поля $(v-v_s)/v_s=\delta/\gamma_s^2+b_2\delta^2+\ldots$ и $b_2=-3\beta_s^2/(2\gamma_s^2)$. Разложение~\eqref{eq:f25-5} удерживает $\Delta v/v_s=\delta/\gamma^2$ и квадрат этого линейного куска; член $b_2$ в него не входит.

Электростатический абзац~\cite{Senichev2013} говорит, что в дефлекторе кинетическая энергия переходит в потенциальную, разброс скорости $\Delta v/v$ уменьшается в $F_v$ раз, траектория при этом почти не меняется, и в~\eqref{eq:f25-6} появляется дополнительный множитель $1/F_v$. Напечатана формула
\begin{align}
\frac{\Delta T_{\mathrm{rev}}}{T_{\mathrm{rev}}}
&=
\left(\alpha_0-\frac{1}{F_v}\frac{1}{\gamma^2}\right)\cdot\delta
\nonumber\\
&\quad
+
\left(
\alpha_1
-
\frac{1}{F_v^2}\frac{\alpha_0}{\gamma^2}
+
\frac{1}{F_v^2}\frac{1}{\gamma^4}
\right)\cdot\delta^2
+
\left(\frac{\Delta L}{L}\right)_{\!\beta}.
\label{eq:f25-25}
\end{align}
Здесь $\gamma$ --- лоренц-фактор синхронной частицы, $\gamma_s$ в обозначениях раздела~\ref{sec:setup}. Линейный член содержит $1/(F_v\gamma^2)$. В квадратичном члене при $\alpha_0/\gamma^2$ и при $1/\gamma^4$ напечатано $F_v^2$, а $(\Delta L/L)_\beta$ стоит с коэффициентом $1$.

Равномерная подстановка «везде $1/\gamma^2\to 1/(F_v\gamma^2)$», которую и описывает фраза про множитель $1/F_v$ в~\eqref{eq:f25-6}, устроена иначе. Квадрат $(\Delta v/v_s)^2$ получает $1/F_v^2$, а перекрёстный член $(\Delta C/C)(\Delta v/v_s)$ получает одну степень $F_v$:
\begin{equation}
\label{eq:f25-logic}
\begin{aligned}
\eta_0^{(F)}
&=
\alpha_0-\frac{1}{F_v\gamma_s^2},
\\
\eta_1^{(F)}
&=
\alpha_1-\frac{\alpha_0}{F_v\gamma_s^2}+\frac{1}{F_v^2\gamma_s^4}.
\end{aligned}
\end{equation}
В~\eqref{eq:f25-25} у $\alpha_0$ во втором порядке стоит $F_v^2$: перекрёстный член в~\eqref{eq:f25-logic} несёт одну степень $F_v$. Ниже с точным рядом сравниваются обе записи. Коэффициенты $\alpha_0$, $\alpha_1$ в обеих --- коэффициенты длины~\eqref{eq:set-C} того кольца, для которого пишется период. Фраза про почти неизменную траекторию значит именно это: длина остаётся в формуле со своими $\alpha_0$ и $\alpha_1$, а $F_v$ стоит при скорости. У сферического дефлектора $\alpha_0=\gamma_s^2+1$ само получено из электрической орбиты; магнитное $\alpha_0$ с последующей поправкой только на $F_v$ его не заменяет.

\subsection{Множитель в тексте IPAC'11}
\label{sec:f25-k}

Формула (25) ссылается на~\cite{Senichev2011} (ссылка [3] работы~\cite{Senichev2013}, спиновые аберрации). Коэффициент, который делит приращение скорости, стоит в соседней работе той же конференции о нелинейном движении в электростатическом кольце~\cite{Zyuzin2011}; символа $F_v$ там нет. Там переход кинетической энергии в потенциальную назван колебательным, потому что тюн кольца не целый, и дальше дословно: процесс конверсии подавлен краевым полем некоторым фактором $k$; этот коэффициент зависит от начальной координаты и убывает с её ростом. После подстановки $k$ в локальную формулу для скорости
\begin{equation}
\label{eq:f25-k}
\frac{\Delta v}{v}
=
\frac{1}{k}
\left[
\frac{\cos\phi-1}{R_{\mathrm{eq}}}\,x_0
+
\sin\phi\cdot x_0'
\right]
+
\frac{\Delta v_0}{v}.
\end{equation}
На $k$ делится только часть, порождённая входными $x_0$ и $x_0'$. Приращение $\Delta v_0/v$, с которым частица вошла в дефлектор, остаётся прежним. Коэффициент $k$ определяется краевым полем и начальной координатой.

В~\eqref{eq:f25-25} множитель устроен иначе. $F_v$ --- одно число на кольцо: им уменьшены члены с $\delta$, а бетатронное удлинение пути стоит с коэффициентом $1$. Коэффициент $k$ из~\eqref{eq:f25-k} уменьшает координатный вклад в $\Delta v$ внутри дефлектора и оставляет $\Delta v_0$ как есть. Это другое место в формуле: $k$ --- локальная характеристика дефлектора и начальной координаты, а не одно число на кольцо.

\subsection{Какое $F_v$ воспроизводит линейный член}
\label{sec:f25-lin}

Отклонение $\delta$ определено в~\eqref{eq:set-delta} полной энергией при нулевом потенциале. Вне поля такая частица имеет скорость
\begin{equation}
\label{eq:f25-vout}
\frac{v-v_s}{v_s}
=
\frac{\delta}{\gamma_s^2}+O(\delta^2).
\end{equation}
На замкнутой орбите часть той же полной энергии сидит в $q\Phi$, и скорость другая. Линейный расчёт раздела~\ref{sec:es} для гладкого электрического кольца даёт
\begin{equation}
\label{eq:f25-eta0}
\eta_0
=
\alpha_0\left(1+\frac{1}{\gamma_s^2}\right)-\frac{1}{\gamma_s^2}
=
\alpha_0-\frac{1-\alpha_0}{\gamma_s^2}.
\end{equation}
Это частный случай $\eta_0=\alpha_0+\alpha_0^{E}/\gamma_s^2-1/\gamma_s^2$ при $\alpha_0^{E}=\alpha_0$: весь изгиб электрический, и линейный потенциал набирается на той же доле орбиты, что и удлинение пути. Сравнивая с первым порядком $\tau=(C/C_s)(v_s/v)-1$ и с~\eqref{eq:set-C}, получаем скорость на орбите
\begin{equation}
\label{eq:f25-vorb}
\frac{v-v_s}{v_s}
=
\frac{1-\alpha_0}{\gamma_s^2}\,\delta+O(\delta^2).
\end{equation}
Отношение внеполевого приращения~\eqref{eq:f25-vout} к орбитальному~\eqref{eq:f25-vorb} равно
\begin{equation}
\label{eq:f25-fv}
F_v
=
\frac{1}{1-\alpha_0}.
\end{equation}
При таком $F_v$ линейная строка~\eqref{eq:f25-logic} совпадает с~\eqref{eq:f25-eta0} целиком. Другого значения, которое сохраняло бы и $\alpha_0$ длины, и этот $\eta_0$, нет: $F_v$ фиксируется отношением двух скоростей. В кольце со смешанным изгибом то же рассуждение с~\eqref{eq:es-eta0} даёт $F_v=1/(1-\alpha_0^E)$.

Смысл $F_v$ прямой. Это число, которое говорит, во сколько раз изменение скорости на смещённой орбите меньше изменения скорости той же полной энергии вне поля. Потенциал забирает долю энергии, и до орбиты доходит множитель $1-\alpha_0$. При $\alpha_0<1$ скорость на внешней орбите растёт медленнее, чем вне поля, и $F_v>1$. При $\alpha_0>1$ она меняется в другую сторону, и $F_v<0$.

Два гладких дефлектора с полем $E\propto r^{-k}$ проверяют~\eqref{eq:f25-fv} до конца, потому что для них $\alpha_0$, $\alpha_1$, $\eta_0$ и $\eta_1$ считаются из одного баланса силы и энергии, без отдельной решётки. Показатель $k=1$ --- цилиндрический дефлектор, $k=2$ --- сферический дефлектор в медианной плоскости. Линейный член потенциала на опорной орбите при любом $k$ один и тот же: поле там фиксировано равновесием синхронной частицы, $q\Phi=\beta_s^2\gamma_s mc^2\,u+O(u^2)$, $u=(r-\rho)/\rho$. Показатель $k$ входит в то, как кинетическая энергия зависит от радиуса.

У цилиндрического дефлектора $\gamma_{\mathrm{kin}}-1/\gamma_{\mathrm{kin}}$ не зависит от радиуса, поэтому на каждой круговой орбите кинетическая энергия равна опорной и скорость равна $v_s$. Весь избыток полной энергии, заданный $\delta$, уходит в потенциал, радиус сдвигается как $u=\delta$, и
\begin{equation}
\label{eq:f25-k1}
\alpha_0=1,
\quad
\alpha_1=\frac{\gamma_s^2+1}{2\gamma_s^2},
\quad
\eta_0=1,
\quad
\eta_1=\frac{\gamma_s^2+1}{2\gamma_s^2}.
\end{equation}
Здесь $F_v$ расходится: орбитальное $\Delta v/v$ в~\eqref{eq:f25-vorb} равно нулю, а $\eta_0^{(F)}$ при $F_v\to\infty$ равно $\alpha_0=1$. Граница напечатанного условия $F_v>1/(\alpha_0\gamma^2)$ при~\eqref{eq:f25-fv} и $F_v>0$ есть $\gamma_s^2=1/\alpha_0-1$, то есть нуль~\eqref{eq:f25-eta0}. Отдельного механизма, помимо этого нуля, линейный член не содержит. У гладкого цилиндрического дефлектора $\alpha_0=1$, этот нуль уходит в $\gamma_s=0$, и $\eta_0=1$ при любой энергии.

У сферического дефлектора кинетическая энергия на внешней орбите падает. Чтобы набрать тот же избыток полной энергии, радиус сдвигается сильнее:
\begin{equation}
\label{eq:f25-k2}
\begin{aligned}
\alpha_0&=\gamma_s^2+1,
&
\eta_0&=\gamma_s^2+2,
\\
\alpha_1&=3\gamma_s^2+\frac{3}{2}+\frac{1}{2\gamma_s^2},
&
\eta_1&=\frac{9\gamma_s^4+7\gamma_s^2+2}{2\gamma_s^2}.
\end{aligned}
\end{equation}
Формула~\eqref{eq:f25-fv} даёт $F_v=-1/\gamma_s^2$. Скорость на орбите в линейном порядке есть $(v-v_s)/v_s=-\delta$: при положительном $\delta$ она падает, а вне поля при том же $\delta$ выросла бы на $\delta/\gamma_s^2$. Линейная строка~\eqref{eq:f25-logic} при этом $F_v$ равна $\gamma_s^2+2$ и совпадает с $\eta_0$ из~\eqref{eq:f25-k2}. При отрицательном $F_v$ напечатанное неравенство $F_v>1/(\alpha_0\gamma^2)$ уже не отмечает знак $\eta_0$: у сферического дефлектора оно не выполнено, а $\eta_0=\gamma_s^2+2$ положителен при любой энергии и нуля не имеет.

\subsection{Квадратичный член}
\label{sec:f25-quad}

Подстановка $\Delta v/v_s=\delta/(F_v\gamma_s^2)$ в~\eqref{eq:f25-5} даёт $\eta_1^{(F)}$ из~\eqref{eq:f25-logic} и по-прежнему отбрасывает $\delta^2$ в самой скорости. Точный период на тех же орбитах удерживает $(v-v_s)/v_s=b_1\delta+b_2\delta^2$. Если $b_1=\delta$-коэффициент~\eqref{eq:f25-vorb}, разложение~\eqref{eq:f25-5} с полными $b_1$ и $b_2$ воспроизводит и $\eta_0$, и $\eta_1$. Отброшенный $b_2$ переходит в период со знаком минус:
\begin{equation}
\label{eq:f25-gap}
\eta_1-\eta_1^{(F)}=-b_2.
\end{equation}
Это тот же пропуск, что $3\beta_s^2/(2\gamma_s^2)$ в магнитной скобке: там $-b_2$ свободной скорости равно $3\beta_s^2/(2\gamma_s^2)$. В дефлекторе в~\eqref{eq:f25-gap} входит $b_2$ скорости на орбите.

У цилиндрического дефлектора скорость на круговых орбитах равна $v_s$ при любом радиусе, поэтому $b_1=b_2=0$ и $F_v\to\infty$. Оба члена $\eta_1^{(F)}$ с $F_v$ в знаменателе исчезают, $\eta_1^{(F)}=\alpha_1$, и это точное $\eta_1$ из~\eqref{eq:f25-k1}. Напечатанные степени $F_v^2$ в~\eqref{eq:f25-25} при $F_v\to\infty$ дают то же самое. Для гладкого цилиндрического дефлектора формула (25) по $\delta$ и по $\delta^2$ верна: задержка равна удлинению пути, потому что скорость вдоль семейства орбит не меняется.

У сферического дефлектора $b_1=-1$ и
\begin{equation}
\label{eq:f25-b2}
b_2=-\frac{\gamma_s^4+1}{2\gamma_s^2},
\qquad
\eta_1-\eta_1^{(F)}=\frac{\gamma_s^4+1}{2\gamma_s^2}.
\end{equation}
При $F_v=-1/\gamma_s^2$ логика~\eqref{eq:f25-logic} даёт
\begin{equation}
\label{eq:f25-sph-logic}
\eta_1^{(F)}=\frac{8\gamma_s^4+7\gamma_s^2+1}{2\gamma_s^2},
\end{equation}
а точный коэффициент --- $\eta_1$ из~\eqref{eq:f25-k2}. Разность~\eqref{eq:f25-b2} не равна $3\beta_s^2/(2\gamma_s^2)$ и не равна тому же выражению, делённому на $F_v$ или на $F_v^2$. Одного числа $F_v$, уже фиксированного линейным отношением скоростей, на второй порядок не хватает: у скорости на орбите есть собственный член $\delta^2$.

Напечатанная степень $F_v^2$ при $\alpha_0$ меняет квадратичный коэффициент ещё раз. При том же $F_v=-1/\gamma_s^2$ формула~\eqref{eq:f25-25} даёт
\begin{equation}
\label{eq:f25-sph-print}
\eta_1^{\mathrm{(25)}}
=
-\gamma_s^4+2\gamma_s^2+\frac{5}{2}+\frac{1}{2\gamma_s^2}.
\end{equation}
С~\eqref{eq:f25-k2} это выражение не совпадает. Сводка для $F_v=1/(1-\alpha_0)$ приведена в таблице~\ref{tab:f25}.

\begin{table}[htb]
\caption{Коэффициенты периода гладкого кольца при $F_v=1/(1-\alpha_0)$. Столбец «логика» --- $\Delta v/v_s=\delta/(F_v\gamma_s^2)$ в разложении~\eqref{eq:f25-5}. Столбец «как напечатано» --- степени $F_v$ из~\eqref{eq:f25-25}.}
\label{tab:f25}
\begin{ruledtabular}
\begin{tabular}{llccc}
дефлектор & & точный & логика & как напечатано \\
\hline
цилиндрический & $\eta_0$ & $1$ & $1$ & $1$ \\
$k=1$, $F_v\to\infty$ & $\eta_1$ & $\dfrac{\gamma_s^2+1}{2\gamma_s^2}$ & то же & то же \\
\hline
сферический & $\eta_0$ & $\gamma_s^2+2$ & $\gamma_s^2+2$ & $\gamma_s^2+2$ \\
$k=2$, $F_v=-1/\gamma_s^2$ & $\eta_1$ & $\dfrac{9\gamma_s^4+7\gamma_s^2+2}{2\gamma_s^2}$ & $\dfrac{8\gamma_s^4+7\gamma_s^2+1}{2\gamma_s^2}$ & \eqref{eq:f25-sph-print} \\
\end{tabular}
\end{ruledtabular}
\end{table}

\subsection{Что в формуле верно}
\label{sec:f25-what}

Линейный член~\eqref{eq:f25-25} верен при $F_v$ из~\eqref{eq:f25-fv}: это $\alpha_0$ длины минус орбитальная скорость~\eqref{eq:f25-vorb}. Для гладкого кольца он совпадает с~\eqref{eq:f25-eta0} и у цилиндрического, и у сферического дефлектора. Квадрат $(\Delta v/v_s)^2$ в~\eqref{eq:f25-5} тоже на своём месте, пока в $\Delta v/v_s$ подставлен линейный кусок. Его не хватает как полного $\eta_1$: рядом с ним должен стоять $-b_2$, аналог магнитного $3\beta_s^2/(2\gamma_s^2)$. У цилиндрического дефлектора этот аналог равен нулю, и второй порядок~\eqref{eq:f25-25} верен. У сферического он равен $(\gamma_s^4+1)/(2\gamma_s^2)$ и в формулу не вошёл.

Степень $F_v^2$ при $\alpha_0/\gamma^2$ в~\eqref{eq:f25-25} не следует из множителя $1/F_v$ в~\eqref{eq:f25-6}. Перекрёстный член длины и скорости несёт одну степень $F_v$, как в~\eqref{eq:f25-logic}. На цилиндрическом дефлекторе расхождение степеней не видно, потому что $F_v$ бесконечен. На сферическом напечатанная степень даёт~\eqref{eq:f25-sph-print} и расходится и с точным $\eta_1$, и с~\eqref{eq:f25-sph-logic}.

Бетатронное слагаемое в~\eqref{eq:f25-25} перенесено из~\eqref{eq:f25-6} с коэффициентом $1$. Это коэффициент удлинения пути. На смещённой бетатронной траектории скорость тоже меняется вместе с потенциалом, и в период входит ещё и этот вклад. Его формула (25) не содержит; как он собирается в $\Lambda$ ряда~\eqref{eq:set-tau}, показано в разделе~\ref{sec:es}.

Формула (25) --- запись периода, аналог~\eqref{eq:f25-6}, а не формула равновесного сдвига. Префактора $1/\eta_0$ в ней нет. В формуле (10) той же работы такой префактор напечатан:
\begin{equation}
\label{eq:f25-10}
\Delta\delta_{\mathrm{eq}}
=
\frac{\gamma_s^2}{\gamma_s^2\alpha_0-1}
\left[
\frac{\delta_m^2}{2}
\left(\alpha_1-\frac{\alpha_0}{\gamma_s^2}+\frac{1}{\gamma_s^4}\right)
+
\left(\frac{\Delta L}{L}\right)_{\!\beta}
\right],
\end{equation}
и $\gamma_s^2/(\gamma_s^2\alpha_0-1)=1/\eta_0$ при магнитном $\eta_0=\alpha_0-1/\gamma_s^2$. Знак плюс. Уравнение второго порядка для $\delta$ в той же работе стоит так, что среднее от квадратичного члена и от $(\Delta L/L)_\beta$ компенсируется линейным членом с минусом. То же минус следует из~\eqref{eq:set-tau}: среднее $\tau$ по синхротронному колебанию обращается в нуль при
\begin{equation}
\label{eq:f25-eq}
\delta_{\mathrm{eq}}
=
-\frac{1}{\eta_0}
\left(
\frac{\eta_1\delta_m^2}{2}
+
\Lambda
\right),
\end{equation}
где $\langle\delta^2\rangle=\delta_m^2/2$ для первой гармоники с амплитудой $\delta_m$, как в определении~\eqref{eq:set-deq}. Плюс в~\eqref{eq:f25-10} этому условию противоречит. В~\eqref{eq:f25-25} той знаковой ошибки нет, потому что равновесный сдвиг там не выписан. Если собрать его из~\eqref{eq:f25-25} тем же шагом, каким из~\eqref{eq:f25-6} получена~\eqref{eq:f25-10}, плюс появится снова; в самой (25) его нет. Знаменатель при этом надо брать из линейного члена (25), $\eta_0=\alpha_0-1/(F_v\gamma_s^2)$, а в числителе стоит тот $\eta_1$, который реально входит в период, вместе с минусом~\eqref{eq:f25-eq}.
